% --- Fate's Edge SRD — Patrons (SRD Cut) ---
\clearpage
\section{Patrons, Runes, and Invokers}
\label{sec:patrons-srd}

\begin{quote}
\emph{“You wish to walk the road of power? Then listen well. The world is old, and older still are the voices beneath it. We call them Patrons, though they were never sworn to us. They are the tides that move unseen, the keepers of forgotten bargains, the sleepers beneath the stone and the stars. To call upon them is to dip a hand into a river that has carved mountains.”}
\end{quote}

\subsection{What Are Patrons?}
Patrons are vast intelligences – not gods, though some worship them as such. They embody concepts and forces: freedom, ruin, shadow, law, the sea, the wild hunt. Each offers power, but always with a cost: fatigue, scars upon fate, or a slow unweaving of one’s own story. To entreat a Patron is to risk being marked. Their Rites are gifts and snares both.

The following ten Patrons are suitable for most campaigns. They are chosen for **ease of use, utility, build diversity, and lore presence**. GMs may add others from the full rulebook.

% ============================================================
\subsection{The Patrons}
% ============================================================

\subsubsection{Oath of Flame \& Light}
\emph{Domain: Dawn, sacred vows, radiant power.} \\
The Oath is no patron of half‑measures. Its fire names, binds, and burns – demanding that those who swear within its radiance stand openly and speak truly. \textit{Gift: Command.}

\subsubsection{Ikasha, She Who Sleeps}
\emph{Domain: Shadow, thresholds, latent potential.} \\
Ikasha is the hush between footfalls, the patience of dark water. In stillness she gathers what might be. \textit{Gift: Stealth.}

\subsubsection{Old Man of the Forest (Grimmir)}
\emph{Domain: Nature, seasons, wilderness.} \\
The Old Man walks the threshold between planted row and forest edge. He speaks in root and stone, fur and feather – the treaty between civilisation and the wild. \textit{Gift: Survival.}

\subsubsection{The Traveler, Lord of the Open Road}
\emph{Domain: Ways, journeys, safe passage.} \\
The Traveler is not merely one who shows the path – they are the path, existing in the pause between steps and in the choice of which fork to take. \textit{Gift: Notice.}

\subsubsection{The Witness, Keeper of Uncomfortable Truths}
\emph{Domain: Truth, revelation, memory.} \\
The Witness remembers what others bury. Every shadow cast and oath broken is a line in her unending ledger. \textit{Gift: Notice.}

\subsubsection{Sealed Gate, Keeper of Thresholds}
\emph{Domain: Boundaries, containment, wards.} \\
The Sealed Gate stands where realities meet, embodying the sacred power of thresholds. Its followers are wardens and exorcists. \textit{Gift: Tinker.}

\subsubsection{Raéyn, Mistress of the Sea}
\emph{Domain: Tides, change, storms, sea travel.} \\
Raéyn is the tempestuous goddess of the sea, mother to all who sail. Her voice is the wind that fills sails; her moods are the storms that test every mariner. \textit{Gift: Skirmish.}

\subsubsection{Maelstraeus, the Infernal Bargainer}
\emph{Domain: Commerce, exchange, contracts.} \\
Maelstraeus is the shadow at every crossroads of exchange. He does not trade fairly – he trades inevitably. \textit{Gift: Persuade.}

\subsubsection{Carrion King, Lord of Renewal}
\emph{Domain: Decay, renewal, transformation.} \\
The Carrion King is not lord of rot for its own sake – he is the alchemist of endings, seeing every corpse as compost for future growth. \textit{Gift: Survival.}

\subsubsection{Mykkiel, Arbiter of the Covenant}
\emph{Domain: Law, zeal, covenant, judgment.} \\
Mykkiel is the unwavering judge of covenants, the divine lawgiver whose scales and sword uphold sacred order. \textit{Gift: Command.}

% ============================================================
\subsection{Patron’s Gift (Imbuement)}
\label{subsec:patron-gift-srd}
Any character with the \textbf{Familiar} (Thiasos) talent may invoke their Patron’s Gift once per scene. Touch an item to imbue it until the end of the scene with \textbf{+1 Melee} and \textbf{+1 Thematic Skill} (see table below). \emph{Push It:} extend one more scene by marking \textbf{+1 Obligation}. Gifts from the same Patron do not stack.

\begin{table}[h]
\centering
\renewcommand{\arraystretch}{1.1}
\begin{tabular}{|p{4cm}|p{3.5cm}|p{7cm}|}
\hline
\textbf{Patron} & \textbf{+1 Thematic Skill} & \textbf{Lore / Gift Bestowal} \\
\hline
Oath of Flame \& Light & Command & Grants the fire of dawn, a vow that shields the faithful and sears the faithless. \\
\hline
Ikasha & Stealth & Grants the hush between footsteps and the raven’s omen at every threshold. \\
\hline
Old Man of the Forest (Grimmir) & Survival & Grants the wild memory: fang, fire, and the path of instinct through the dark wood. \\
\hline
The Traveler & Notice & Grants the open way, a compass that never rests, and roads where none are marked. \\
\hline
The Witness & Notice & Grants the unblinking gaze that unmasks deceit and remembers every oath. \\
\hline
Sealed Gate & Tinker & Grants mastery of thresholds – doors that yield or bar at your command. \\
\hline
Raéyn & Skirmish & Grants the sailor’s fortune: winds that shift, storms that answer to will. \\
\hline
Maelstraeus & Persuade & Grants the contract’s weight, every deal sealed in fire and shadow. \\
\hline
Carrion King & Survival & Grants the feast of decay, where what is dead becomes seed for what lives. \\
\hline
Mykkiel & Command & Grants the blazing verdict, a covenant carried in fire and oath, binding law to blade. \\
\hline
\end{tabular}
\caption{Patron’s Gift: Thematic Skill and lore. Bonus applies only when fiction matches the Patron’s domain.}
\label{tab:patron-gift-srd}
\end{table}

\subsection{Terrestrial Patrons (Mortal Factions)}
Not all patrons are cosmic. Mortal factions offer power without supernatural debt – but always with strings.

\begin{itemize}
  \item \textbf{Guild (Merchant / Craft)} – Perks: access to markets, resources, skilled labor, legal loopholes. Obligation: pay dues, protect members, enforce charter.
  \item \textbf{Noble House} – Perks: political influence, safe houses, armed retainers. Obligation: attend court, back house feuds, marry strategically.
  \item \textbf{Crime Syndicate} – Perks: black market, fences, smuggled goods, protection rackets. Obligation: pay tribute, silence witnesses, eliminate rivals.
  \item \textbf{Military Order} – Perks: elite training, weapons, rank, tactical support. Obligation: follow chain of command, accept deployments, obey martial law.
  \item \textbf{Temple / Faith} – Perks: sanctuary, healing, relics, pilgrim networks. Obligation: tithe, spread doctrine, perform rituals.
\end{itemize}

\textbf{Using Terrestrial Patrons}
\begin{itemize}
  \item Treat them as a single Obligation track (capacity = Spirit + Presence).
  \item Calling on a perk adds +1 Obligation (or more for major requests).
  \item When Obligation fills, the patron makes a demand – refuse and lose all perks and gain a Rival faction.
  \item Clear Obligation by performing services for the patron (downtime).
\end{itemize}

\subsection{Rites, Obligation, and the Invoker}
Characters who swear to a single Patron (Runekeepers) learn specific \textbf{Rites} – magical invocations purchased with XP. Invoking a Rite takes 1 action and marks \textbf{+1 Obligation}. \emph{Push It:} mark an additional +1 Obligation to amplify the effect.

Invokers do not fully swear; they carry \textbf{Patron’s Symbols}. Invoking a Rite through ritual takes DV+1 rounds and always marks +1 Obligation. Alternatively, they may \emph{Crack the Seal} – cast instantly by setting the Symbol to \textsc{Compromised} and marking +2 Obligation (+3 for High Rites).

Obligation capacity = Spirit + Presence. Exceeding capacity marks 1 Fatigue per segment over. If ever over double capacity, clear all Fatigue, mark +1 Harm, and the Patron intrudes (GM frames a thematic demand). Obligation can be reduced by performing acts of service aligned with the Patron’s domain (downtime activity, subject to GM approval).

\begin{quote}
\emph{“Every Rune is a promise. Every line a covenant. Do not mistake the Runekeeper’s silence for weakness; their memory is the foundation of our craft.”}
\end{quote}